\paragraph{Automated MS lesion segmentation.}
Current MS lesion segmentation tools have improved mainly through network
design, ensembling, and training data, while keeping voxel-level training
objectives. Examples include LST-AI, an ensemble of three 3D U-Nets trained
with Tversky and cross-entropy loss~\cite{wiltgen_lst-ai_2024}, the
nnU-Net-based FLAMeS~\cite{dereskewicz_flames_2025}, and SegHeD+, which
trains on heterogeneous datasets with longitudinal and anatomical
constraints~\cite{basaran_seghed_2024}; the ICPR 2024 challenge on MSLesSeg
benchmarked such methods across hospitals and
timepoints~\cite{rondinella_icpr_2024}. Although these tools generalize
well across centers, small lesions remain a major source of error: most
lesions missed by FLAMeS were smaller than
10~mm$^3$~\cite{dereskewicz_flames_2025}. Lesion-wise detection metrics
have therefore been advocated over the Dice score because they reflect what
matters for clinical reading~\cite{basit_rethinking_2026}, and
ConfLUNet segments individual lesions, including confluent
ones~\cite{wynen_conflunet_2025}. These works change how lesions are
evaluated or separated at inference, but not how the model is trained.

\paragraph{Training objectives for small lesions.}
Alternatives to the standard objective reweight voxels by class volume,
error type, or difficulty~\cite{sudre_generalised_2017,salehi_tversky_2017,lin_focal_2017,abraham_novel_2019},
or balance lesions of different sizes by computing the loss per
lesion~\cite{shirokikh_universal_2020,kofler_blob_2023,bouteille_learning_2026,bouteille_bicc_2026,molchanova_benchmarking_2025}.
All of them still score each lesion by how well it is outlined, not by
whether it is found. Losses that target detection directly model lesions as
fixed shapes~\cite{zhang_all-net_2021,rachmadi_improving_2024}, and multiple
instance learning, which treats an object as found if any of its parts is
found~\cite{dietterich_solving_1997,ilse_attention-based_2018}, has been
used in segmentation almost only when voxel-level labels are
missing~\cite{pinheiro_image-level_2015,cheng_fmamil_2025}. CATMIL
(\cref{sec:method}) adds to the standard nnU-Net
objective~\cite{isensee_nnu-net_2021} the question of whether each
annotated lesion was detected.
